\documentclass{beamer}

\usepackage{fontspec}
\usepackage{listings}
\usepackage{xspace}
\usepackage[backend=biber]{biblatex}

\setmainfont{TeX Gyre Heros}
\setmonofont{Fantasque Sans Mono}
\lstset{basicstyle=\ttfamily\footnotesize}
\usetheme{focus}
\parskip=1\baselineskip
\def\lstlanguagefiles{lstlean.tex}
\lstset{language=lean}

\definecolor{keywordcolor}{rgb}{0.7, 0.1, 0.1}   % red
\definecolor{commentcolor}{rgb}{0.4, 0.4, 0.4}   % grey
\definecolor{symbolcolor}{rgb}{0.4, 0.4, 0.4}    % grey
\definecolor{sortcolor}{rgb}{0.1, 0.5, 0.1}      % green

\newcommand{\N}{\mathbb{N}}
\newcommand{\Q}{\mathbb{Q}}
\newcommand{\lean}[1]{\texttt{#1}\xspace} % for writing Lean expressions
\newcommand{\ex}{\lean{ex}}
\newcommand{\mathlib}{\texttt{mathlib}\xspace}
\newcommand{\pow}{{\ensuremath{{}^\wedge}}\xspace}
\newcommand{\ring}{\lean{ring}}
\newcommand{\ringexp}{\lean{ring\_exp}}

\addbibresource{lean.bib}

\title{A Lean tactic for normalising ring expressions with exponents}
\author{Anne Baanen}
\titlegraphic{\includegraphics[scale=.25]{lf_vu.png}}
%\titlegraphic{\includegraphics[scale=.2]{vu.png}}
\institute{\vspace{-\baselineskip}Lean Forward\\Vrije Universiteit Amsterdam}
\date{\vspace{-\baselineskip}IJCAR, 2 July 2020}

\begin{document}

\frame{\titlepage}

\begin{frame} \frametitle{Background}
	Lean is a proof assistant based on the calculus of constructions~\cite{lean-prover}.\\
	It has a simple kernel for proof checking and an elaborator with powerful tactic support.\\
	The Lean community is developing \texttt{mathlib},
	a repository of formalised classical mathematics proofs and proof automation.

	The Lean Forward project aims to make proof assistants\\
	accessible to mathematicians by developing proof automation informed by users' needs.
\end{frame}

\begin{frame} \frametitle{Background}
	To a mathematician, the following is obvious:
	\begin{align*}
		2^{n+1} - 1 = 2 * 2^n - 1
	\end{align*}
	Lean should do it automatically.

	The previous \ring tactic~\cite{ring-tactic} uses Horner normal form: efficient for the semiring operators $+$ and $*$,
	but it doesn't support exponentiation $\pow$.

\pause
	It's not too hard to solve this manually:
	rewrite $a^{n + 1} = a^n * a$ and \ring can finish by applying commutativity.

	Such rules don't work unconditionally: $x^{100}$ should not become 100 multiplications of $x$.
\end{frame}

\begin{frame} \frametitle{Design overview}
	\textbf{Goal:} a practical normalising tactic for expressions with $+$, $*$ and \pow, numerals (in $\Q$) and variables.
	It should solve all goals that \ring can and be approximately as fast.

\pause
	To prove $a = b$, normalise $a$ giving $p_a : a = a'$ and normalise $b$ giving $p_b : b = b'$, then check $a'$ is identical to $b'$.
	If identical, $p_a p_b^{-1}$ proves $a = b$.

	Lean tactics typically construct proof terms directly (no reflection), since the elaborator is faster than the kernel.
\end{frame}

\begin{frame}[fragile]
	The normal form is a syntax tree in the type family \lean{ex}.
	The children for each node are restricted by a parameter \lean{ex\_type}:

\begin{lstlisting}
inductive ex_type : Type
| sum | prod | exp | base
inductive ex : ex_type → Type
| zero  : ex_info →                     ex sum  -- 0
| sum   : ex_info → ex prod →  ex sum → ex sum  -- +
| coeff : ex_info →             coeff → ex prod -- rat
| prod  : ex_info → ex exp  → ex prod → ex prod -- *
| exp   : ex_info → ex base → ex prod → ex exp  -- ^
| var   : ex_info →              atom → ex base -- atom
| sum_b : ex_info →            ex sum → ex base -- (...)
\end{lstlisting}

	\begin{itemize}
		\item $(a + b) + c$ is not allowed: left argument to \lean{sum} must be a product
		\item $a * (b + c)$ is not allowed: right argument to \lean{prod} must be a product
	\end{itemize}
\end{frame}

\begin{frame} \frametitle{Intricacies}
Commutativity: pick a linear order $\prec$ on \lean{ex}.
Then sort $a + b + c + \cdots$ so that $a \prec b \prec c \prec \cdots$.

\pause
To limit expression size, don't unfold $100 * a$ to $a + a + \cdots + a$.\\
This means keeping track of coefficients.\\
The function \lean{add\_overlap} decides when to add coefficients:
\begin{align*}
	&\lean{add\_overlap}\ (3 * x^2)\ (7 * x^2) &&= 10 * x^2 \\
	&\lean{add\_overlap}\ (3 * x^2)\ (7 * y^2) &&= 3 * x^2 + 7 * y^2 \\
	&\lean{add\_overlap}\ (3 * x^2)\ (-3 * x^2) &&= 0 &\text{(not $0 * x^2$)}
\end{align*}
\pause
In a general semiring $R$, exponentiation has type $\pow : R \to \N \to R$.\\
During execution, \ringexp keeps track of the current type\\
using a reader monad transformer.
\end{frame}

\begin{frame} \frametitle{Optimisations}
To be practical, the \ringexp tactic must be fast.\\
The Horner form used by the \ring tactic is optimal for $+$ and $*$.\\
Optimisation is needed to achieve acceptable running time.

\pause
Typeclass instances and implicit arguments cost time to infer,\\
so they are cached as much as possible:\\
instances are stored with the current type in the reader monad,\\
implicit arguments and intermediate values in the \lean{ex\_info} field of \lean{ex}.

\pause
Result: on average a constant factor slowdown compared to \ring.\\
Noticeably faster on problems with larger exponents\\
($20$ times on $x^{50} * x^{50} = x^{100}$),
also in practice for $(1 + x^2 + x^4 + x^6) * (1 + x) = 1 + x + x^2 + x^3 + x^4 + x^5 + x^6 + x^7$.
\end{frame}

\begin{frame}
\printbibliography
\end{frame}
\end{document}
